\chapter{Využití práce a integrace měřicího systému do vozidla}
\label{chap:integrace}

\section{Vazba na celkovou senzoriku formule}

Navrhovaný RTK GNSS/IMU systém tvoří část mého širšího úkolu zajistit
senzoriku studentské formule a připravit data pro následnou analýzu.
Semestrální práce vytváří podklad pro bakalářskou práci: od výběru
komponent přecházím k návrhu společného sběru dat, časového přiřazení
a instalace do skutečného automobilu. Tato kapitola zachycuje
\textbf{stav návrhu k 13.~říjnu~2026}, nikoli hotovou kabeláž,
funkční akviziční program nebo provedené jízdní zkoušky.

Trajektorii chci spojit s otáčkami jednotlivých kol, pohybem mechanismů
roll/heave, teplotami brzdových třmenů a povrchu pneumatik. Získám tak
možnost hledat souvislosti mezi místem na trati a chováním vozidla:
například porovnat průjezdy zatáčkou, odpružení při brzdění nebo zahřívání
pneumatik. Poloha sama tyto jevy nevysvětluje; prostorově přiřazuje další měření.

Po posouzení ceny a předpokládané délky přívodů přibližně \SI{1.5}{m}
měním původní návrh dvou nápravových uzlů na \textbf{centrální sběr u
Raspberry Pi}. Analogové signály povedu jednotlivě do společného ADC
a vstupních obvodů v krytu loggeru. Senzorové CAN 3 a CAN 4 ani jejich
mikrokontroléry nyní nezařazuji do výchozí sestavy. Dva CANy zachovávám
pouze pro příjem dat auta a předávání vybraných údajů ADU. Distribuovaný
sběr zůstává zálohou, pokud centrální varianta nevyhoví při měření.

Základním požadavkem je uchovat v Raspberry Pi co nejúplnější původní
záznam. ADU nemusí přijímat všechny hodnoty ani jejich plnou frekvenci.
Odděluji proto \emph{měření a archivaci} od \emph{informování řidiče}.
Logger v první etapě neřídí brzdy, pohon ani jiné bezpečnostně významné
funkce. Jeho porucha nesmí být podmínkou funkce vozidla; případné využití
dat pro řízení by vyžadovalo samostatný návrh a validaci.

\section{Měřené veličiny a pracovní předpoklady}

Původní označení snímačů jako měření odklonu kola nebylo přesné.
\textbf{Odklon kola v tomto návrhu neměřím.} Sleduji pohyb mechanismů
roll/heave. Úhel hřídele snímače není automaticky úhlem karoserie ani
zdvihem kola; převod odvodím z kinematiky a kalibrace skutečného mechanismu.
Pro dimenzování zatím předpokládám dva polohové kanály vpředu a dva vzadu.
Počet i uchycení potvrdím nad modelem vozidla.

U brzd volím pro první verzi kontaktní měření \textbf{třmenu}, nikoli
bezkontaktní měření kotouče. Výsledek tedy není okamžitou teplotou třecí
plochy ani brzdové kapaliny. Teploty a polohy mají rozdílnou požadovanou
šířku pásma; všem čidlům automaticky nepřiřazuji \SI{500}{Hz}.

\begin{table}[htbp]
\centering\small
\caption{Pracovní rozsah centrálního DAQ nad rámec GNSS/IMU. Frekvence jsou cíle návrhu.}
\label{tab:vozidlo-senzory}
\begin{tabularx}{\textwidth}{@{}L{29mm}L{16mm}L{32mm}Y@{}}
\toprule
Veličina & Počet & Záznam & Cesta k Raspberry Pi \\
\midrule
Poloha roll/heave & 4 kanály & oživení 200 Hz; cíl 500 Hz/kanál & analogový přívod, ochrany, filtr, ADC \\
Otáčky kol & 4 snímače & každá platná hrana & úprava úrovně, chráněné GPIO \\
Teplota třmenů & 4 snímače & 4 Hz/kanál & termočlánek, speciální převodník u loggeru \\
Teplota pneumatik & 4 snímače & cíl 10 Hz/kanál & ověřit analogový výstup nebo převod PWM \\
Data auta & CAN 1 & původní zprávy & externí CAN řadič, listen-only \\
\bottomrule
\end{tabularx}
\end{table}

\subsection{Polohové snímače a kandidát centrálního ADC}

Pro náhradu opotřebitelného potenciometru zvažuji magnetický AS5600
s analogovým výstupem. Jeho 12bitový rozsah odpovídá kroku
$360^{\circ}/4096\approx0{,}088^{\circ}$, nikoli přesnosti mechanismu.
Potřebuje správně ustavený magnet, držák a ochranu proti vodě. Levný
modul není přímou mechanickou náhradou utěsněného potenciometru.
Ověřím také napěťový rozsah, filtraci a odezvu \cite{as5600-integrace}.
Pro tuto architekturu upřednostním analogový výstup před vzdáleným
I\textsuperscript{2}C; stejné adresy čidel pak nejsou problémem.
Kalibraci může usnadnit nastavení elektrického rozsahu na skutečně
využitý úhel. Dva 12bitové převody za sebou nezaručují 12bitovou přesnost.

Raspberry Pi 5 nemá na 40pinovém konektoru analogové měřicí vstupy.
Přívod k Raspberry proto znamená přívod k \emph{centrálnímu převodníku},
nikoli připojení proměnného napětí přímo na GPIO \cite{rpi-gpio-adc}.
Jako levného kandidáta pro první zkoušku volím \textbf{MCP3208}:
osm multiplexovaných vstupů, 12bitový SAR převod a SPI. Výrobce uvádí
nejvýše 100~tisíc převodů/s při \SI{5}{V} a 50~tisíc při
\SI{2.7}{V}; jde o celkový výkon převodníku, ne každého vstupu
\cite{mcp3208-adc}. Prototyp předpokládá logiku \SI{3.3}{V}
a SPI \SI{1}{MHz}. Napájení, referenci a mezní parametry ještě ověřím
pro skutečné zapojení. Při provozu ADC na \SI{5}{V} by bylo třeba
přizpůsobit také logické úrovně SPI, nikoli jen analogové vstupy.

První čtyři kanály vyhradím poloze, další čtyři zůstanou pro teploty
pneumatik, pokud uspěje níže uvedený převod na napětí. Při jiném řešení
IR čidel tvoří rezervu. Rozlišení 12~bitů nepovažuji za definitivně
dostatečnou přesnost: nejprve stanovím požadovanou chybu polohy podle
převodového poměru mechanismu. Pokud nevyhoví rozlišení, šum nebo
časový souběh, porovnám přesnější či simultánně vzorkující ADC.

\subsection{Teploty a otáčky nejsou univerzální analogové vstupy}

Čtyři termočlánky typu K povedu odpovídajícím termočlánkovým či
kompenzačním vedením ke čtyřem převodníkům MAX6675 u loggeru.
Jejich SPI zůstane krátké. Maximální převod \SI{220}{ms} vede k cíli
\SI{4}{Hz}; rozlišení \SI{0.25}{\degreeCelsius} není přesností měření
\cite{max6675-integrace}. Ověřím teplotní rozsah sondy, kompenzaci
studeného konce i elektrické spojení s třmenem. Přechod na měděné vodiče
a vstupní svorky musí být tepelně posouzeny; převodníky nesmějí trpět
teplotním gradientem od Pi nebo měniče. Montáž nezasáhne do hydrauliky
ani bezpečnostně významného spojení brzdy. Termočlánek nepřipojím jako
běžné napěťové čidlo přímo k MCP3208.

Pro pneumatiky zůstává kandidátem MLX90614, například BAA. Široké zorné
pole varianty xAA je nominálně $90^{\circ}$, proto ověřím výhled při
propružení a řízení. Čidlo nabízí SMBus a programovatelný PWM výstup,
nikoli hotový analogový výstup \cite{mlx90614-integrace}. Pro centrální
variantu chci na stole vyzkoušet \textbf{PWM přívod a centrální filtr
s převodem střídy na napětí}. Jde o můj návrh dalšího obvodu, nikoli
ověřenou vlastnost zakoupeného modulu. Je nutné kalibrovat úrovně,
rozsah, zbytkové zvlnění a zpoždění, rozpoznat přerušení vodiče a
neplatný výstup. Samotný ADC kód nebude vydáván za původní digitální
teplotu čidla. Pokud cesta nevyhoví, porovnám vhodné analogové IR čidlo
nebo upravené digitální rozhraní; tuto část nákupu zatím neuzavírám.

Cíl čtení IR kanálu je \SI{10}{Hz}; skutečná čerstvost údajů závisí
na čidle a filtrech. Jeden snímač na kolo neměří profil celého běhounu.
Ověřím emisivitu, znečištění a propustnost případného okénka.
Bez ověření kapacity vedení a časování nevolím běžný SMBus/I\textsuperscript{2}C
po celé délce kabelu. Přípustnost takového spoje nelze určit pouze
z délky \SI{1.5}{m} \cite{nxp-i2c-adc}.

Otáčky kol chci snímat digitálním bezkontaktním čidlem. Konkrétní typ
ještě neuzavírám: Hallovo čidlo s magnetem a indukční čidlo s kovovým
cílem mají jiné požadavky na terč. Každý přívod zakončím u loggeru
ochranou a úpravou úrovně, podle signálu také hysterezním vstupem.
Výstup \SI{5}{V} nebo \SI{12}{V} nesmí přímo na GPIO. Časované
události hran v Linuxu mají být výchozím řešením, nikoli smyčka
opakovaně dotazující úroveň pinu. Prokáže-li test ztráty nebo
nepřijatelnou časovou chybu, doplním \emph{centrální} hardwarový čítač
či zachytávání hran; není nutné kvůli tomu ihned vracet uzly ke kolům.

\section{Centrální zapojení, kabeláž a rozhraní}

\subsection{Proč nyní dávám přednost centrálnímu sběru}

Původně jsem uvažoval dva místní mikrokontroléry a dvě senzorové CAN
větve. Při omezeném rozpočtu a přívodech přibližně \SI{1.5}{m} chci
nejprve ověřit jednodušší instalaci s jednou vstupní deskou u loggeru.
Odpadnou dva programované uzly, jejich zdroje, kryty a samostatná
synchronizace hodin. Neodpadnou ochrany ani převody potřebné pro
jednotlivé typy čidel. Cenová úspora proti distribuované variantě musí
zahrnovat také větší kabelový svazek a případné přesnější analogové vstupy.

\begin{table}[htbp]
\centering\small
\caption{Rozhraní výchozí centralizované sestavy; nikoli finální pinová mapa.}
\label{tab:centralni-rozhrani}
\begin{tabularx}{\textwidth}{@{}L{34mm}L{43mm}Y@{}}
\toprule
Zdroj & Dlouhý úsek k loggeru & U Raspberry Pi \\
\midrule
Poloha roll/heave & analogový signál s měřicí zemí & ochrana, filtr, ADC a krátké SPI \\
Otáčky kol & impulzy s definovanou referencí & úprava úrovně a GPIO události \\
Termočlánky třmenů & termočlánkové vedení & 4 speciální převodníky, krátké SPI \\
IR teplota & pracovní varianta PWM & ověřit převod na napětí a ADC \\
GNSS a IMU & podle navigační části & USB, SPI/INT a GNSS PPS \\
Palubní elektronika & CAN 1 a CAN 2 & dva externí CAN kanály \\
\bottomrule
\end{tabularx}
\end{table}

Centrální desku umístím na odpruženou část společně s loggerem. Přívody
rozdělím konektorově na přední a zadní svazek, ale nepřejmenuji tyto
analogové kabely na CAN 3 a CAN 4. Skutečné délky, pohyblivé úseky,
ohyby a servisní přístup ověřím v CAD. Délka \SI{1.5}{m} je zatím
pracovní předpoklad, ne doložená délka všech vodičů.

\subsection{Analogová kvalita a ochrana vstupů}

Signál snímače povedu společně s jeho měřicí referencí a oddělím jej
od vysokoproudových a spínaných vedení. Senzorová zem nesmí být
náhodným bodem karoserie, přes který teče výkonový proud. Zapojení
stínění určím podle skutečné zemní topologie a výsledků měření.
U každého vstupu navrhnu omezení poruchového proudu, přepěťovou ochranu
a analogový filtr. Ten musí potlačit i složky, které by při vzorkování
způsobily aliasing; digitální filtr již vzniklý aliasing neodstraní.

Výstup například \SIrange{0.5}{4.5}{V} je nutné přizpůsobit zvolenému
rozsahu ADC. Dělič, tolerance reference, teplotní drift a pokles
napájení snímače vstupují do nejistoty celé sestavy. Poměrové měření
má smysl jen při skutečně známé vazbě napájení snímače a reference ADC.
Příliš vysoká výstupní impedance nebo nevhodný RC filtr mohou zhoršit
ustálení multiplexovaného převodu; případný buffer navrhnu podle
vzorkovacího vstupu ADC, nikoli jen podle statického napětí.

Rozdíl zemních potenciálů může v jednovodičovém napěťovém měření
vypadat jako změna polohy. Pokud tato chyba nebude přijatelná,
porovnám diferenciální vstupní zesilovač nebo vhodnější ADC
\cite{adi-vstupy-adc}. Pseudodiferenciální režim MCP3208 není
libovolně plovoucí měření: záporný vstup smí být podle datasheetu jen
v blízkosti země, přibližně \(\pm\SI{100}{mV}\)
\cite{mcp3208-adc}. Neřeší automaticky všechny zemní rozdíly v autě.

Před návrhem celé desky připojím jeden skutečný snímač přes skutečný
kabel. Porovnám šum a posun nuly s krátkým přívodem, při vypnutém
a zapnutém pohonu, ventilátorech a ostatních spotřebičích. Podmínky
testu včetně reference, napájení a filtru uložím. Za dostatečnou
nepovažuji pouze vizuálně hladkou křivku po silném vyhlazení.

\subsection{CAN 1 a CAN 2 zůstávají rozhraním auta}

\begin{table}[htbp]
\centering\small
\caption{Platné označení CAN rozhraní po přechodu na centrální ADC.}
\label{tab:can-role}
\begin{tabularx}{\textwidth}{@{}L{17mm}L{32mm}L{25mm}Y@{}}
\toprule
Název & Hlavní tok dat & Rychlost & Účel \\
\midrule
CAN 1 & auto $\rightarrow$ RPi & 1 Mbit/s & potvrzená rychlost; listen-only \\
CAN 2 & RPi $\rightarrow$ auto & asi 500 kbit/s & ověřit; vybrané údaje pro ADU \\
CAN 3/4 & nyní neosazeny & neurčena & pouze označení dřívější záložní varianty \\
\bottomrule
\end{tabularx}
\end{table}

Šipky vyjadřují hlavní aplikační tok, ne elektrickou jednosměrnost CAN.
Na CAN 1 chci pouze pasivní příjem bez vysílání a potvrzování.
Na CAN 2 připustím jen dohodnuté identifikátory a frekvence, nikoli
automatické přeposílání celého provozu. U zpráv určím měřítko, jednotky,
pořadí bajtů, znaménko, neplatné hodnoty a čas platnosti. Poznámku
o 16bitových kanálech ADU ověřím podle jeho přesného modelu; není
to omezení samotné CAN sběrnice.

Pi 5 potřebuje externí CAN řadiče. Pro dvě rozhraní nyní stačí
uvažovat \emph{jeden} Waveshare 2-CH CAN HAT+ s podporou Pi 5
\cite{waveshare-can-integrace}. Alternativou jsou samostatné řadiče
nebo vhodný USB--CAN. MCP2515 má jen dvě přijímací vyrovnávací paměti
\cite{mcp2515-integrace}; úplnost příjmu CAN 1 musím ověřit i při
souběžném ADC, GNSS, IMU a zápisu SSD. Nízké vytížení CPU samo
neprokazuje bezeztrátový záznam. Rozhraní zpřístupním přes SocketCAN
\cite{socketcan-integrace} a stabilně přiřadím význam CAN 1/2.

Zakončení existujících větví nejprve zjistím. Správný nový CAN má dvě
zakončení na fyzických koncích, nikoli odpor u každého zařízení;
připojení loggeru proto neznamená přidání dalšího odporu
\cite{ti-can-integrace}. Zachovám potřebné galvanické oddělení.
Před objednávkou sestavím společnou pinovou mapu ADC, CAN, IMU SPI/INT,
PPS, čtyř vstupů otáček a teplotních převodníků. Nestačí mechanická
kompatibilita HATů; nesmějí kolidovat volby čipů, přerušení ani
napájení. Prověřím také M.2 HAT+, chladič a přístup ke konektorům.

\section{Výpočet vzorkování, přenosu a objemu dat}
\label{sec:can-vypocet}

\subsection{Propustnost ADC a místního SPI}

Uvažované polohové kanály sdílejí převodník, proto požaduji
\begin{equation}
 f_{\mathrm{ADC,celkem}}=\sum_i f_i.
\end{equation}
Čtyři kanály při \SI{500}{Hz} znamenají 2~000 jednotlivých převodů/s.
Pro osm stejně rychlých kanálů je to 4~000 převodů/s. Jde o výkon
čipu; neprokazuje pravidelnost akvizičního programu pod Linuxem.
Oblíbený ADS1115 má nejvýše 860 převodů/s celkem \cite{ads1115-adc};
jeden kus tedy čtyři kanály po \SI{500}{Hz} nezvládne. Ani teoretických
$860/4=215$ vzorků/s na kanál nemá pro \SI{200}{Hz} rozumnou časovou rezervu.

Pro odhad SPI počítám se třemi přenesenými bajty na jeden převod,
tedy 24 hodinovými impulzy a taktem \SI{1}{MHz}. Je to model
konkrétního způsobu čtení, nikoli univerzální režie všech ADC.
Pro podíl obsazení hodin SPI pak platí
\begin{equation}
 U_{\mathrm{SPI}}=\frac{24\,f_{\mathrm{ADC,celkem}}}{10^6}.
\end{equation}

\begin{table}[htbp]
\centering\small
\caption{Výpočtová zátěž centrálního ADC. Režie OS a ostatních zařízení není zahrnuta.}
\label{tab:adc-vypocet}
\begin{tabularx}{\textwidth}{@{}YR{26mm}R{28mm}R{20mm}@{}}
\toprule
Režim & Převody/s & Hodinové bity/s & SPI 1 MHz \\
\midrule
4 polohy po 200 Hz & 800 & 19 200 & 1,92 \% \\
4 polohy po 500 Hz & 2 000 & 48 000 & 4,80 \% \\
4 polohy po 500 Hz + 4 IR po 10 Hz & 2 040 & 48 960 & 4,90 \% \\
8 kanálů po 500 Hz, rezervní model & 4 000 & 96 000 & 9,60 \% \\
\bottomrule
\end{tabularx}
\end{table}

Jedna taková transakce trvá na vodičích přibližně \SI{24}{\micro\second}.
Čtyři po sobě čtené polohy tedy obsadí asi \SI{96}{\micro\second};
nejde o simultánní měření. Do periody \SI{2}{ms} při \SI{500}{Hz}
se musí vejít i přepínání CS, systémová volání, další zařízení a čekání.
Je-li stejný řadič sdílen s IMU a CAN, zahrnu i jejich přenosy a
priority. Nesmím zaměnit maximální takt SPI zařízení s jeho
vzorkovací frekvencí. Zlepšení rozlišení či menší šum získaný filtrací
musím posoudit společně s přidaným zpožděním.

\subsection{Frekvence impulzů kol a záznam na SSD}

Pro rychlost $v=100/3{,}6\ \mathrm{m/s}$, pracovní průměr kola
$D=\SI{0.46}{m}$ a $N=12$ cílů s jednou zaznamenanou hranou na cíl platí
\begin{equation}
 f_{\mathrm{hrany}}=\frac{vN}{\pi D}\approx\SI{230.66}{Hz}
 \quad\text{na jedno kolo}.
\end{equation}
Pro čtyři kola jde přibližně o 922,64 události/s. Rozměr kola,
počet cílů i rychlost jsou \textbf{výpočtové předpoklady}, nikoli
potvrzené parametry auta. Prokluz a vyšší rychlost počet zvýší;
záznam obou hran místo jedné jej zdvojnásobí. Události již neposílám
senzorovým CANem, ale přijímám je na chráněných vstupech loggeru.

Pro hrubý odhad úložiště si stanovuji vlastní binární záznam
\SI{16}{\byte} na hodnotu či hranu: 64bitový čas, 32bitová hodnota,
16bitový kanál a 16bitový stav. Hlavička bloku má navíc pořadí,
verzi formátu a kontrolu integrity. Pole hodnoty pojme původní ADC
kód; fyzikální jednotku určuje uložená konfigurace. Následující
výpočet ještě nezahrnuje blokové hlavičky, CAN ani GNSS/IMU.

\begin{table}[htbp]
\centering\small
\caption{Model velikosti nových senzorových záznamů při 16 B na položku.}
\label{tab:adc-uloziste}
\begin{tabularx}{\textwidth}{@{}YR{26mm}R{31mm}@{}}
\toprule
Zdroj a pracovní režim & Položky/s & MB/h, desetinně \\
\midrule
4 polohy po 500 Hz & 2 000 & 115,20 \\
4 kola, uvedený příklad & 922,64 & 53,14 \\
4 třmeny po 4 Hz & 16 & 0,92 \\
4 pneumatiky po 10 Hz & 40 & 2,30 \\
\midrule
Celkem, rychlý režim & 2 978,64 & 171,57 \\
Celkem při poloze 200 Hz & 1 778,64 & 102,45 \\
\bottomrule
\end{tabularx}
\end{table}

Součty zaokrouhluji až po výpočtu. Změna balení více kanálů do bloku
může režii snížit; komprese nesmí maskovat ztracené vzorky. Datový
objem není hlavní obtíží, důležitější je pravidelnost čtení,
neblokující ukládání a dohledatelnost výpadků.

\subsection{Zatížení palubních CAN sběrnic}

Původní výpočty zatížení dvou senzorových CAN větví nejsou po změně
architektury rozpočtem instalovaného systému. CAN 1 závisí na skutečném
provozu auta; nejprve pořídím pasivní záznam a zjistím četnosti,
délky rámců, špičky i ztráty v přijímači.

Pro CAN 2 uvažuji čtyři osmibajtové zprávy po \SI{10}{Hz}. U klasického
CAN s 11bitovým identifikátorem používám konzervativní délku
135 bitových časů včetně bit stuffingu a mezisnímkového intervalu
\cite{bosch-can-integrace}. Nezahrnuje chyby ani opakování přenosu.
Při předběžné rychlosti \SI{500}{kbit\per\second} vychází
\begin{equation}
 \Delta U_{\mathrm{CAN\,2}}=
 \frac{4\cdot10\cdot135}{500\,000}=1{,}08\,\%.
\end{equation}
Jde pouze o přírůstek, nikoli celkové vytížení CAN 2. Obsahem mohou
být stav loggeru, platnost navigace a vybrané teploty. Vysokofrekvenční
polohy a původní hrany kol ukládám místně; jejich plný tok ADU nepotřebuje.

\section{Časová osa, kalibrace a struktura záznamu}

Centralizace odstraní potřebu synchronizovat dva vzdálené MCU,
neodstraní však nejistotu okamžiku měření. MCP3208 převádí během
transakce SPI a nemá vlastní volně běžící vzorkovací plán s FIFO.
Čas návratu funkce v uživatelském programu není přesný čas vzorku.
V první implementaci zaznamenám monotónní časy ohraničující přenos,
pořadí kanálů a naměřenou nejistotu. Skutečný okamžik vzorkování
ověřím na CS/SCLK logickým analyzátorem. Pravidelná smyčka se zpožděním
\SI{2}{ms} sama nezaručuje rovnoměrných \SI{500}{Hz}.

Pro otáčky použiji GPIO character-device API a události z jádra
namísto časování až při probuzení uživatelského programu. Linux
poskytuje pořadová čísla událostí, podle nichž lze odhalit přetečení
fronty \cite{linux-gpio-adc}. To ale nezaručuje zachycení každého
fyzického impulzu: nesmí být kratší než možnosti vstupu a jeho
obsluhy. Hardware generující známý počet impulzů proto použiji jako
nezávislou kontrolu. Neprohlašuji předem, že Pi 5 poskytne pro každý
takový vstup přesný hardwarový timestamp hrany.

U rychlých kanálů si stanovuji pracovní cíl časového přiřazení do
\SI{1}{ms}. Ověřím jitter, maximální mezeru mezi vzorky i časový rozdíl
kanálů při zatížení všemi periferiemi. Pro roll/heave může být
mezikanálový posun významný i při správné průměrné frekvenci.
Nesplnění cíle může vést nejprve k ADC se spouštěním či FIFO nebo
k centrálnímu zachytávání hran, teprve potom k distribuované variantě.

Vazbu monotónního času loggeru na GNSS poskytne PPS a časové zprávy
popsané v kapitole~\ref{chap:vyber}. Transformaci na GNSS čas uložím
odděleně, aby změna systémových hodin nepřepsala historii. Při
\SI{100}{km\per\hour} odpovídá časová chyba \SI{1}{ms} posunu
\SI{2.78}{cm}; nejde o úplnou chybu polohy ani o již dosaženou přesnost.
Teplotní převody a filtry mají vlastní odezvu, kterou musím popsat
vedle časového razítka čtení.

Původní tři datové proudy GNSS, IMU a online navigace zůstávají.
Přidám \texttt{analog.bin} s ADC kódy, \texttt{wheel\_edges.bin}
s hranami, \texttt{temperature.bin} s odečty teplotních převodníků
a interpretací IR kanálů a oddělené záznamy CAN 1/2. Kalibrované
hodnoty nesmějí přepsat původní odečty. Manifest jízdy uchová
reference, převodní rozsahy, filtry, kalibrace, mapu konektorů,
verze programu a DBC. U IR převodu přes filtr je původním záznamem
pro logger ADC kód, nikoli neuchované digitální slovo uvnitř čidla.

Záznam rozdělím do kratších uzavíraných bloků s kontrolou integrity.
Při odpojení čidla nebo vynechání převodu uložím stav a mezeru,
nikoli opakovanou starou hodnotu označenou jako nový vzorek.

\section{Napájení z palubní sítě 8\textendash16~V}
\label{sec:napajeni-vozidlo}

\subsection{Výkonový návrh a kandidát měniče}

Vozidlo má elektronické jištění přes palubní PMU. Napájení DAQ
navrhuji z vyhrazeného jištěného výstupu. Měnič musí pracovat při
\SIrange{8}{16}{V}, jmenovitě \SI{12}{V}; PMU jej nenahrazuje.
Pro větev \SI{5}{V} nejprve rezervuji \SI{25}{W} a zdroj posuzuji
i při \SI{30}{W}. Nejde o změřenou spotřebu. Do bilance nově patří
Pi, SSD, GNSS, jeden dvoukanálový CAN HAT, ventilátor, centrální ADC,
vstupní obvody a napájené snímače. Samostatné nápravové MCU nyní
neuvažuji. Odběr, rozběh a rezervu změřím v celé sestavě; snížení
počtu CAN HATů neznamená automaticky známý pokles spotřeby.

Dosavadním referenčním kandidátem je Pololu D42V55F5, číslo 5571:
výstup \SI{5}{V} s tolerancí \(\pm3\,\%\), vstup až \SI{60}{V}
a typický proud \SI{6}{A} při podmínkách výrobce
\cite{pololu-d42-integrace}. Minimální vstup závisí na úbytku,
trvalý proud na napětí a chlazení. Kvůli ceně ponechávám výběr
konkrétního měniče otevřený; přechod na centrální ADC není schválením
nákupu Pololu ani jiné náhrady. Levnější modul musí projít stejným
zátěžovým testem, nejen mít příznivý údaj na štítku.

Při pracovním předpokladu $\eta=0{,}90$ platí
\begin{equation}
 I_{\mathrm{in}}=\frac{P_{\mathrm{out}}}{\eta U_{\mathrm{in}}}.
\end{equation}
Pro \SI{30}{W} je to \SI{4.17}{A} při \SI{8}{V}, \SI{2.78}{A}
při \SI{12}{V} a \SI{2.08}{A} při \SI{16}{V}; ztráta by byla
\SI{3.33}{W}. Účinnost je předpoklad, nikoli změřený parametr.
Je nutné ověřit i teplotu měniče v uzavřeném krytu.

Předchozí cenový podklad pro Pololu byl 870~Kč s DPH k
12.~říjnu~2026 \cite{botland-pololu-integrace}. Dřívější plánovací
částka přibližně 1~200~Kč s DPH na zdroj a základní připojení
není závaznou cenou ani schváleným výběrem. Neobsahuje zálohování.
Původní rozpočet přibližně 24~290~Kč bez DPH nadále popisuje
\textbf{GNSS/IMU elektroniku, nikoli kompletní DAQ vozidla}.

\subsection{Ochrany, přívod a vypínání}

Rozsah \SIrange{8}{16}{V} neznamená odolnost proti všem přechodovým
jevům. Ani vysoké maximální napětí samo neprokazuje automobilovou
kvalifikaci; na omezení testování upozorňuje i Pololu
\cite{pololu-auto-integrace}. Navrhnu vstupní filtr a přepěťovou
ochranu podle skutečné palubní sítě. TVS, energie impulzu a svorkové
napětí zatím nejsou uzavřeným schématem. Spínaný zdroj zároveň
oddělím a odfiltruji od citlivé ADC reference a měřicích zemí.

Samostatnou tavnou pojistku nepřidávám automaticky, protože jištění
má zajistit PMU. Její nastavení však musí chránit i nejslabší vodič
a konektor. Slabší senzorové odbočky vyžadují samostatné posouzení.
Proud \SI{4.17}{A} na vstupu při \SI{8}{V} nelze zaměnit s proudem
na \SI{5}{V} výstupu. Podle měření určím běžný limit, rozběhový proud,
jeho dobu a reakci na poruchu. Dokumentace Ecumaster tyto veličiny
rozlišuje \cite{ecumaster-pmu-integrace}; model, firmware a kanál
PMU ve vozidle ještě doplním. Po dobu záznamu požaduji spojité
napájení, nikoli PWM napájení měniče. PMU není zárukou čistého
výstupu \SI{5}{V}, stabilní analogové reference ani funkcí UPS.

Měnič umístím blízko Pi; po autě povedu \SI{12}{V}, ne dlouhý
vysokoproudový přívod \SI{5}{V}. Úbytky změřím přímo na Pi i čidlech.
Při napájení přes GPIO použiji dimenzované pevné propojení s oběma
\SI{5}{V} piny a odpovídajícími zemními kontakty. Laboratorní
propojky nejsou finální kabeláž. Zdroje na USB-C, GPIO nebo HATu
nesmějí být neřízeně spojeny paralelně.

Bez USB Power Delivery může být nutné nastavení
\texttt{PSU\_MAX\_CURRENT=5000}. Nezlepší měnič ani vodiče;
použiji ji až po ověření celé cesty \cite{rpi-power-integrace}.
Pokles pod \SI{8}{V} současný požadavek nepokrývá. Snižující měnič
není UPS a náhlé vypnutí může ztratit data ve frontě. Ověřím obnovu
uzavřených bloků a s týmem vyřeším řádné vypnutí nebo samostatně
navržený doběh. Neobejdu přitom bezpečnostní vypínání vozidla.

\section{Instalace, nákupní rozsah a využití dat}

Centrální logger se vstupní deskou umístím na odpruženou část do
servisně přístupného krytu chráněného proti vodě a nečistotám.
V CAD prověřím obálku Pi, SSD, HATu, převodníků a konektorů,
ohybové poloměry i chlazení. Pohyblivé svazky odlehčím v tahu
a ověřím celý zdvih i řízení. Nesmějí zasáhnout do kola nebo brzdy.
3D tištěný kryt sám neprokazuje tepelnou ani vibrační odolnost.

IMU bude tuze a opakovatelně ustavena, GNSS anténa bude mít výhled
a vhodnou montáž. Zaměřím jejich vzájemný vektor a polohu vůči
referenčnímu bodu auta. Měnič, výkonové vodiče a magnety polohových
snímačů rozmístím s ohledem na magnetometr. Kalibraci provedu až
v sestavě a zkontroluji při zapnutých spotřebičích.

\begin{table}[htbp]
\centering\small
\caption{Změna nákupního rozsahu proti distribuovanému návrhu. Položky DAQ ještě nejsou úplně oceněny.}
\label{tab:adc-nakup}
\begin{tabularx}{\textwidth}{@{}L{38mm}L{26mm}Y@{}}
\toprule
Celek & Pracovní počet & Stav po změně \\
\midrule
Centrální ADC & 1 & kandidát MCP3208, doplnit referenci a vstupní obvody \\
Polohové vstupy & 4 & ochrany, přizpůsobení, filtrace a případné buffery \\
Vstupy otáček & 4 & ochrany a úprava úrovně pro GPIO \\
Termočlánkové převodníky & 4 & centrální umístění, kandidát MAX6675 \\
IR vstupní cesta & 4 & ověřit PWM převod; není hotový nákup \\
CAN rozhraní u Pi & 2 kanály & kandidát jeden 2-CH CAN HAT+ \\
Nápravové MCU a CAN 3/4 & 0 & pouze záložní alternativa, nyní neobjednávat \\
Kabeláž a napájení & podle CAD & senzory, konektory, kryt a zdroj ocenit samostatně \\
\bottomrule
\end{tabularx}
\end{table}

Původní BOM a krátký nákupní seznam GNSS/IMU tímto neměním na údajný
úplný rozpočet auta. Centrální deska ještě vyžaduje schéma, skutečný
seznam součástek a cenové porovnání. Vynechané uzly tedy nepřevádím
na smyšlenou přesnou úsporu. Upřednostním české či evropské dodavatele
na fakturu a sloučení objednávek; možnost sponzorského dodání je
alternativou, nikoli jistým nulovým nákladem.

Pro roll/heave vytvořím kalibrační převod podle geometrie, ne pouze
přejmenování kanálů. Rychlost kol odvodím z časů hran a valivého obvodu.
Porovnání s GNSS rychlostí může ukázat prokluz, ale musí zohlednit
zatáčení, rozdílné rychlosti kol a polohu antény. Při malé rychlosti
je poměrné vyjádření prokluzu špatně podmíněné. Odvozené veličiny
nesmějí nahradit původní záznam.

Teploty přiřadím místu na trati s vědomím tepelné setrvačnosti čidel
a třmenu. Mezi jízdami zachovám souřadnou soustavu a ustavení základny.
Uložím nastavení podvozku, pneumatik a kalibrací, aby změna sestavy
nebyla zaměněna za chybu měření.

\section{Postup ověření a návrat k distribuované variantě}

\begin{enumerate}
 \item \textbf{Inventura.} Potvrdím počet a rozsah snímačů, výstupní
       napětí, model ADU, rychlost CAN 2, zakončení a PMU. V CAD určím
       délky a konektory; připravím společnou pinovou a napájecí mapu.
 \item \textbf{Jeden analogový kanál.} Se skutečným kabelem změřím
       chybu proti referenci, šum, ustálení, dynamickou odezvu a odpojení.
       Ověřím rozsah a ochrany nejprve na stole, nikoli poruchou v autě.
 \item \textbf{Celý centrální vstup.} Zprovozním čtyři polohy nejprve
       na 200 a potom na 500 Hz/kanál. Přidám generátor impulzů kol,
       termočlánky a ověřovací IR cestu. Změřím jitter a posun kanálů;
       u teplot rozliším četnost čtení a čerstvost měření.
 \item \textbf{Souběh loggeru.} Test poběží s GNSS, IMU, CAN 1/2 a SSD.
       Známým počtem impulzů a rámců ověřím ztráty, fronty a reakci na
       odpojení. PPS a testovací signál ověří společnou časovou osu.
 \item \textbf{Napájení a vozidlo.} Provedu starty a zátěžové testy
       při 8, 12 a 16 V, měření úbytků, teplot a rozběhu přes PMU.
       V autě bez jízdy ověřím rušení, pohyb svazků a CAN režimy.
       Krátká jízda následuje až po kontrole týmem.
\end{enumerate}

Za minimální podmínky první verze považuji alespoň třicetiminutový
souvislý záznam, dohledatelné výpadky, nulový rozdíl mezi vyslanými
a přijatými testovacími událostmi a žádné vysílání na CAN 1.
U rychlých kanálů ověřím cíl časové vazby do \SI{1}{ms} a samostatně
uvedu největší mezeru i mezikanálový posun. Přijatelnou chybu polohy
stanovím až podle mechaniky v odpovídajících jednotkách. Prázdný
chybový log sám není důkazem, že nedošlo ke ztrátě. Po restartu musí
být čitelné uzavřené bloky a označitelná chybějící koncová část.

Pokud nevyhoví analogový přenos, nejprve ověřím kabeláž, reference,
vstupní filtr a případný diferenciální přijímač. Pokud nevyhoví
časování, posoudím centrální hardwarové vzorkování a čítače.
Teprve při prokázaném přínosu vrátím místní měřicí desky s USB,
RS-485 nebo CAN. Při návratu ke CAN mohou oba uzly sdílet jedinou
senzorovou větev; původní rozdělení na CAN 3 a CAN 4 není podmínkou
funkce. Taková změna bude vyžadovat nový výpočet přenosu a synchronizace.

Výsledkem této etapy má být doložený centrální DAQ propojující
bakalářskou práci se senzorikou formule: původní data, společná časová
osa, reprodukovatelné převody a popsaná omezení. ADC, vstupní schéma,
IR cesta, ochrany, mechanika a úplný rozpočet zůstávají podmíněny
ověřením. Návrh centrálního sběru není tvrzením, že už byly tyto
zkoušky úspěšně provedeny.
